%% ARITHMETIQUE %%
%
\section{Rapport et polyn\^ome de $1^{er}$ degr\'e }

\itemstyle{ Op\'erateur~: } \operfmt{Poly1\_XsurY\_flt}
\index{Poly1\_XsurY\_flt}
\label{OP-Poly1_XsurY_flt}

%%%%%OBJET%%%%%
\itemstyle{ Objet~: }
Application d'un polyn\^ome de $1^{er}$ degr\'e 
au rapport de 2 param\`etres.\\
Sp\'ecification en entr\'ee de l'unit\'e, des noms fran\c cais et
anglais du vecteur de sortie.

%%%%%ARGUMENTS%%%%%
\bigskip
%%%%%ENTREES%%%%
\begin{tabarg}
\multicolumn{3}{l}{\titretab {Entr\'ees}}\\
\hline
$X$ & Vecteur, \texttt{float} & Premier op�rande (unit� quelconque)\\
$Y$ & Vecteur, \texttt{float} & Second op�rande (unit� quelconque)\\
$C_0$ & Coeff., \texttt{float} & Coefficient constant \\
$C_1$ & Coeff., \texttt{float} & Coefficient de $X/Y$ \\
$Z_{indef}$ & Coeff., \texttt{float} & Valeur de sortie dans le cas o\`u
le rapport $X/Y$ est ind\'efini \\ 
\textit{Unit\_res} & Coeff., \texttt{String} & Unit� du r�sultat. \\
\textit{Name}   & Coeff., \texttt{string} & Descriptif (�nom long�) du r\'esultat \\
\textit{Eng\_name} & Coeff., \texttt{string} & Descriptif (�nom long�) du r\'esultat en langue anglaise \\
\end{tabarg}

\medskip
%%%%%SORTIES%%%%
\begin{tabarg}
\multicolumn{3}{l}{\titretab {Sorties}}\\
\hline
$Z$ & Vecteur, \texttt{float} & R�sultat (unit� indiqu�e par \textit{Unit\_res})\\
\end{tabarg}


%%%%%FORMULE%%%%
\itemstyle{Formule}

%%%%
\[
Z = C_0 + C_1.\frac{X}{Y}
\]
si $Y=0$  alors $Z = Z_{indef}$


%%%%%COMMENTAIRE%%%%
%\itemstyle{ Commentaire~:}
%
%Pas de traitement sp\'ecifique du cas o\`u un \'echantillon de $Y$ est nul.

%%%%%AUTEUR%%%%%
\itemstyle{Auteur~:}
CNRM/GMEI/TRAMM

%%%%%REFERENCES%%%%%
\itemstyle{R\'ef\'erences~:}
